\begin{derivation}[远区角通量积分与总辐射功率]

从远区辐射场
\[
\vb{E}_{\rm rad}
=
\frac{q_{\mathrm e}}{4\pi\varepsilon_0cR}
\frac{
\vb{n}\times[
(\vb{n}-\vb{\beta}_{\rm rel})
\times\dot{\vb{\beta}}_{\rm rel}]
}
{\mathfrak d_{\rm ret}^{\,3}},
\qquad
\mathfrak d_{\rm ret}
=
1-\vb{n}\cdot\vb{\beta}_{\rm rel}
\]
出发。远区场横向于传播方向，因此
\[
\vb{n}\cdot\vb{E}_{\rm rad}=0,
\qquad
\vb{B}_{\rm rad}
=
\frac1c\vb{n}\times\vb{E}_{\rm rad}.
\]
Poynting 向量的径向分量于是为
\[
\vb{n}\cdot\vb{S}_{\rm rad}
=
\frac1{\mu_0}
\vb{n}\cdot
(\vb{E}_{\rm rad}\times\vb{B}_{\rm rad})
=
\varepsilon_0c|\vb{E}_{\rm rad}|^2.
\]
乘球面元 \(R^2\dd\Omega\)，按\emph{观察时间} \(t\) 计的角能流为
\[
\frac{\dd P_{\rm obs}}{\dd\Omega}
=
\frac{q_{\mathrm e}^2}{16\pi^2\varepsilon_0c}
\frac{
\left|
\vb{n}\times[
(\vb{n}-\vb{\beta}_{\rm rel})
\times\dot{\vb{\beta}}_{\rm rel}]
\right|^2
}
{\mathfrak d_{\rm ret}^{\,6}}.
\]

这里还不能直接把它称为“源在单位时间内发射的功率”，因为不同观察方向对应的推迟时间间隔不同。由
\[
t-t_{\rm r}
=
\frac{R(t_{\rm r})}{c}
\]
求导，并用
\[
\frac{\dd R}{\dd t_{\rm r}}
=
-\vb{n}\cdot\vb{v}
\]
得到
\[
\frac{\dd t}{\dd t_{\rm r}}
=
1-\vb{n}\cdot\vb{\beta}_{\rm rel}
=
\mathfrak d_{\rm ret}.
\]
因此同一辐射能量满足
\[
\frac{\dd E}{\dd t_{\rm r}\dd\Omega}
=
\frac{\dd E}{\dd t\dd\Omega}
\frac{\dd t}{\dd t_{\rm r}},
\]
从而
\[
{
\frac{\dd P_{\rm emit}}{\dd\Omega}
=
\frac{q_{\mathrm e}^2}{16\pi^2\varepsilon_0c}
\frac{
\left|
\vb{n}\times[
(\vb{n}-\vb{\beta}_{\rm rel})
\times\dot{\vb{\beta}}_{\rm rel}]
\right|^2
}
{(1-\vb{n}\cdot\vb{\beta}_{\rm rel})^5}.
}
\]
这就是正文\eqnrefs{eq:lienard_angular_emission_power}。分母从六次幂降到五次幂，正是推迟时间雅可比因子的结果。

总辐射功率可由这个角分布直接积分得到。由于总功率只依赖 \(\vb{\beta}_{\rm rel}\) 与 \(\dot{\vb{\beta}}_{\rm rel}\) 的相对方向，可取
\[
\vb{\beta}_{\rm rel}
=
\beta_{\rm rel}\vb{e}_z,
\qquad
\dot{\vb{\beta}}_{\rm rel}
=
\dot\beta_\perp\vb{e}_x
+
\dot\beta_\parallel\vb{e}_z,
\]
并写
\[
u=\cos\theta,
\qquad
\dd\Omega=\dd\varphi\,\dd u.
\]
对方位角先积分。直接展开双重叉积并使用
\[
\int_0^{2\pi}\cos\varphi\,\dd\varphi=0,
\qquad
\int_0^{2\pi}\cos^2\varphi\,\dd\varphi=\pi
\]
可得
\begin{align*}
&\int_0^{2\pi}\dd\varphi\,
\left|
\vb{n}\times[
(\vb{n}-\vb{\beta}_{\rm rel})
\times\dot{\vb{\beta}}_{\rm rel}]
\right|^2\\
&\qquad=
\pi\Bigl[
2\dot\beta_\parallel^2(1-u^2)
+
\dot\beta_\perp^2
\bigl((1+\beta_{\rm rel}^2)(1+u^2)-4\beta_{\rm rel}u\bigr)
\Bigr].
\end{align*}
平行、垂直两部分因此分别需要两个一维积分：
\[
I_\parallel
=
2\pi
\int_{-1}^{1}
\frac{1-u^2}{(1-\beta_{\rm rel}u)^5}\,\dd u,
\]
\[
I_\perp
=
\pi
\int_{-1}^{1}
\frac{
(1+\beta_{\rm rel}^2)(1+u^2)-4\beta_{\rm rel}u
}
{(1-\beta_{\rm rel}u)^5}\,\dd u.
\]
这些是有理函数积分。令
\[
w=1-\beta_{\rm rel}u,
\qquad
u=\frac{1-w}{\beta_{\rm rel}},
\qquad
\dd u=-\frac{\dd w}{\beta_{\rm rel}},
\]
上下限变为 \(w=1+\beta_{\rm rel}\) 与
\(w=1-\beta_{\rm rel}\)。逐项积分 \(w^{-3},w^{-4},w^{-5}\) 后得到
\[
I_\parallel
=
\frac{8\pi}{3}
\frac1{(1-\beta_{\rm rel}^2)^3}
=
\frac{8\pi}{3}\gamma_{\rm L}^6,
\]
以及
\[
I_\perp
=
\frac{8\pi}{3}
\frac1{(1-\beta_{\rm rel}^2)^2}
=
\frac{8\pi}{3}\gamma_{\rm L}^4.
\]
代回角功率，
\begin{align*}
P_{\rm rad}
&=
\frac{q_{\mathrm e}^2}{16\pi^2\varepsilon_0c}
\left(
I_\parallel\dot\beta_\parallel^2
+
I_\perp\dot\beta_\perp^2
\right)\\
&=
\frac{q_{\mathrm e}^2}{6\pi\varepsilon_0c}
\left(
\gamma_{\rm L}^6\dot\beta_\parallel^2
+
\gamma_{\rm L}^4\dot\beta_\perp^2
\right).
\end{align*}
由于
\[
\dot{\vb{\beta}}_{\rm rel}
=
\frac{\vb{a}}{c},
\qquad
|\vb{\beta}_{\rm rel}\times\vb{a}|^2
=
\beta_{\rm rel}^2a_\perp^2,
\]
并且
\[
\gamma_{\rm L}^6
\left[
a^2-
|\vb{\beta}_{\rm rel}\times\vb{a}|^2
\right]
=
\gamma_{\rm L}^6a_\parallel^2
+
\gamma_{\rm L}^4a_\perp^2,
\]
最终得到
\[
{
P_{\rm rad}
=
\frac{q_{\mathrm e}^2\gamma_{\rm L}^6}
{6\pi\varepsilon_0c^3}
\left[
a^2-
|\vb{\beta}_{\rm rel}\times\vb{a}|^2
\right],
}
\]
即正文\eqnrefs{eq:lienard_total_power_from_angular_flux}。非相对论极限
\(\beta_{\rm rel}\to0\) 自动退化为 Larmor 功率
\[
P_{\rm rad}
\longrightarrow
\frac{q_{\mathrm e}^2a^2}{6\pi\varepsilon_0c^3}.
\]
\end{derivation}

\begin{derivation}[超相对论辐射的形成长度]
真空远场中 $\kappa_z=(\omega/c)\cos\theta$，而 $v=\beta_{\rm rel}c$。因此正文\eqnrefs{eq:formation_length_general_mismatch} 精确给出
\begin{equation*}
L_f
=\frac{c}{\omega\,[\beta_{\rm rel}^{-1}-\cos\theta]}.
\end{equation*}
令 $x=\gamma_{\rm L}^{-2}$，则
$\beta_{\rm rel}^{-1}=(1-x)^{-1/2}$。对函数
$f(x)=(1-x)^{-1/2}$ 在 $x=0$ 处作一阶 Taylor 展开。由于
\begin{equation*}
f''(u)=\frac{3}{4}(1-u)^{-5/2},
\end{equation*}
Lagrange 余项给出某个 $\xi\in(0,x)$ 使
\begin{equation*}
r_\beta(x)
=f(x)-1-\frac{x}{2}
=\frac{f''(\xi)}{2}x^2.
\end{equation*}
因 $f''(u)$ 在 $0\le u<x<1$ 上单调递增，
\begin{equation*}
0\le r_\beta(x)
\le\frac{3x^2}{8(1-x)^{5/2}}.
\end{equation*}

角度部分使用四阶 Taylor 公式。存在介于 $0$ 与 $\theta$ 之间的中间角 $\theta_*$，使
\begin{equation*}
\cos\theta
=1-\frac{\theta^2}{2}
+\frac{\theta^4}{24}\cos\theta_*.
\end{equation*}
所以
\begin{equation*}
r_\theta(\theta)
=1-\cos\theta-\frac{\theta^2}{2}
=-\frac{\theta^4}{24}\cos\theta_*,
\qquad
|r_\theta|\le\frac{\theta^4}{24}.
\end{equation*}
由此
\begin{align*}
\beta_{\rm rel}^{-1}-\cos\theta
&=\frac{x+\theta^2}{2}+r_\beta+r_\theta\\
&=\frac{x+\theta^2}{2}(1+\delta_f),
\end{align*}
其中
\begin{equation*}
\delta_f
=\frac{2(r_\beta+r_\theta)}{x+\theta^2}.
\end{equation*}
三角不等式给出
\begin{equation*}
|\delta_f|
\le
{
\epsilon_f
\equiv
\frac{
\dfrac{3x^2}{4(1-x)^{5/2}}
+\dfrac{\theta^4}{12}
}{x+\theta^2}}.
\end{equation*}
代回形成长度即得到正文\eqnrefs{eq:formation_length_ultrarelativistic}。若 $\epsilon_f<1$，则
\begin{align*}
\left|\frac{L_f-L_f^{(0)}}{L_f^{(0)}}\right|
&=\left|\frac1{1+\delta_f}-1\right|\\
&=\frac{|\delta_f|}{|1+\delta_f|}\\
&\le\frac{\epsilon_f}{1-\epsilon_f},
\end{align*}
其中
$L_f^{(0)}=2c/[\omega(x+\theta^2)]$。这给出了常用超相对论小角形成长度的定量适用条件。
\end{derivation}

\begin{derivation}[源损失中的吸收与辐射通道]
先证明 Green 表达式。对 $\ee^{-\ii\omega t}$ 的时间约定，频域 Maxwell 方程给
\[
\widetilde{\vb{E}}
=\ii\mu_0\omega\,\mathbf G^R\widetilde{\vb{J}}.
\]
定义复数
\[
Z_J
\equiv
\iint\dd^3r\,\dd^3r'\,
\widetilde{\vb{J}}^{*}(\vb{r})\cdot
\mathbf G^R(\vb{r},\vb{r}';\omega)
\cdot\widetilde{\vb{J}}(\vb{r}').
\]
则
\[
\int\dd^3r\,\widetilde{\vb{J}}^{*}\cdot\widetilde{\vb{E}}
=\ii\mu_0\omega Z_J.
\]
电子或外部源损失给场的平均功率取负的机械做功率，
\begin{align*}
P_{\rm loss}
&=-\frac12\Re(\ii\mu_0\omega Z_J)\\
&=\frac{\mu_0\omega}{2}\Im Z_J.
\end{align*}
而
\[
\Im Z_J
=\iint\dd^3r\,\dd^3r'\,
\widetilde{\vb{J}}^{*}(\vb{r})\cdot
\operatorname{Im}\mathbf G^R(\vb{r},\vb{r}';\omega)
\cdot\widetilde{\vb{J}}(\vb{r}'),
\]
得到正文\eqnrefs{eq:source_total_loss_green_power}。

局域能量守恒进一步把总损耗分解为材料吸收和辐射。取非磁介质
$\vb{B}=\mu_0\vb{H}$、
$\vb{D}=\varepsilon_0\boldsymbol\varepsilon_r\vb{E}$。频域 Maxwell 方程为
\[
\nabla\times\vb{E}=\ii\omega\mu_0\vb{H},
\qquad
\nabla\times\vb{H}=\vb{J}-\ii\omega\varepsilon_0
\boldsymbol\varepsilon_r\vb{E}.
\]
第二式取复共轭：
\[
\nabla\times\vb{H}^*
=\vb{J}^*+\ii\omega\varepsilon_0
\boldsymbol\varepsilon_r^\dagger\vb{E}^*.
\]
使用向量恒等式
\[
\nabla\cdot(\vb{E}\times\vb{H}^*)
=\vb{H}^*\cdot(\nabla\times\vb{E})
-\vb{E}\cdot(\nabla\times\vb{H}^*)
\]
得到
\begin{align*}
\nabla\cdot(\vb{E}\times\vb{H}^*)
&=\ii\omega\mu_0|\vb{H}|^2
-\vb{E}\cdot\vb{J}^*\\
&\quad
-\ii\omega\varepsilon_0
\vb{E}\cdot\boldsymbol\varepsilon_r^\dagger\vb{E}^*.
\end{align*}
对体积 $V$ 积分并取实部。第一项纯虚，不贡献平均能流；利用
\[
\Re\!\left[-\ii\,
\vb{E}\cdot\boldsymbol\varepsilon_r^\dagger\vb{E}^*\right]
=-\vb{E}^*\cdot\boldsymbol\varepsilon_I\cdot\vb{E},
\]
以及散度定理，得到
\begin{align*}
\frac12\Re\oint_{\partial V}\dd\vb{A}\cdot
(\vb{E}\times\vb{H}^*)
&=-\frac12\Re\int_V\dd^3r\,\vb{J}^*\cdot\vb{E}\\
&\quad
-\frac{\omega\varepsilon_0}{2}
\int_V\dd^3r\,
\vb{E}^*\cdot\boldsymbol\varepsilon_I\cdot\vb{E}.
\end{align*}
识别
\[
P_{\rm rad}
=\frac12\Re\oint\dd\vb{A}\cdot(\vb{E}\times\vb{H}^*),
\qquad
P_{\rm abs}
=\frac{\omega\varepsilon_0}{2}\int\dd^3r\,
\vb{E}^*\cdot\boldsymbol\varepsilon_I\cdot\vb{E},
\]
以及
$P_{\rm loss}=-(1/2)\Re\int\vb{J}^*\cdot\vb{E}$，就得到
\[
{P_{\rm loss}=P_{\rm abs}+P_{\rm rad}.}
\]
对被动介质 $\boldsymbol\varepsilon_I\succeq0$，故 $P_{\rm abs}\ge0$；采用出射边界条件时 $P_{\rm rad}\ge0$。
\end{derivation}

\begin{derivation}[各向异性导电薄层的模式极点]
先从面内平移不变情形开始。令
\[
\sigma_{ab}(\vb{\rho},\vb{\rho}';\omega)
=
\sigma_{ab}(\vb{\rho}-\vb{\rho}';\omega).
\]
采用二维 Fourier 变换
\[
f(\vb{\rho})
=
\int\frac{\dd^2q}{(2\pi)^2}
f(\vb{q})\ee^{+\ii\vb{q}\cdot\vb{\rho}},
\]
卷积定理把正文\eqnrefs{eq:general_2d_conductivity_kernel} 化成
\[
K_a(\vb{q},\omega)
=
\sum_b
\sigma_{ab}(\vb{q},\omega)
E_{\parallel,b}(\vb{q},0,\omega).
\]

对固定 $\vb{q}$，两侧体介质中的平面波都满足无源 Maxwell 方程。先看 TM 极化。取 $\vb{q}=q\hat{\vb{x}}$，切向电场为 $E_x$、磁场为 $H_y$。由
\[
\vb{H}
=
\frac1{\mu_0\omega}\vb{k}\times\vb{E}
\]
和
$\vb{k}\cdot\varepsilon_j\vb{E}=0$，对沿 $\pm z$ 传播的 TM 波得到
\[
H_y^{(\pm)}
=
\pm
\frac{\varepsilon_0\varepsilon_j\omega}{k_{zj}}
E_x^{(\pm)}
=
\pm Y_{p,j}E_x^{(\pm)}.
\]
TE 极化同理给
\[
H_x^{(\pm)}
=
\mp
\frac{k_{zj}}{\mu_0\omega}
E_y^{(\pm)}
=
\mp Y_{s,j}E_y^{(\pm)}.
\]
所以 $Y_{p,j}$、$Y_{s,j}$ 正是把切向电场映射到切向磁场的波导纳。

现在寻找无外部入射的束缚/共振表面模。取两侧都从界面向外衰减或满足出射条件的解，并令共同切向场幅为 $E_\parallel$。切向电场连续性已经满足。磁场跳跃条件
\[
\hat{\vb{z}}\times(\vb{H}_2-\vb{H}_1)
=
\vb{K}
=
\boldsymbol\sigma\vb{E}_\parallel
\]
则给出一个齐次线性方程。对各向同性二维层，TM 与 TE 不混合，因此分别得到
\[
\bigl(Y_{p,1}+Y_{p,2}+\sigma_L\bigr)E_p=0,
\]
\[
\bigl(Y_{s,1}+Y_{s,2}+\sigma_T\bigr)E_s=0.
\]
存在非零 $E_p$ 或 $E_s$ 的条件正是正文\eqnrefs{eq:2d_sheet_surface_mode_conditions}。

若
\[
\boldsymbol\sigma
=
\begin{pmatrix}
\sigma_{pp}&\sigma_{ps}\\
\sigma_{sp}&\sigma_{ss}
\end{pmatrix}
\]
在 TE/TM 基底中含非对角元，边界方程变成
\[
\left[
\begin{pmatrix}
Y_{p,1}+Y_{p,2}&0\\
0&Y_{s,1}+Y_{s,2}
\end{pmatrix}
+
\boldsymbol\sigma
\right]
\begin{pmatrix}E_p\\E_s\end{pmatrix}
=0.
\]
非零解存在当且仅当系数矩阵行列式为零，即正文\eqnrefs{eq:2d_sheet_anisotropic_pole_determinant}。

最后说明为什么二维空间色散对电子特别重要。电子在距界面有限冲击参数处的 Fourier 电流包含连续 $q$，且近场部分在
$q>\omega/c$ 时对应倏逝真空波。远场光只在接近光锥的横向通道取样，而 EELS 的 Green 核积分会覆盖这些大 $q$ 倏逝分量；它们通过上述边界矩阵直接读取 $\sigma_L(q,\omega)$、$\sigma_T(q,\omega)$ 的非局域结构。因此二维表面电导是三维双点介电核在严格零厚度极限下的自然接口，而不是人为引入的有效参数。

最后从同一边界方程求反射振幅。对 TE 波，切向电场连续给 $t_s=1+r_s$；磁场跳跃条件给
\[
Y_{s,1}(1-r_s)-Y_{s,2}t_s=\sigma_Tt_s.
\]
消去 $t_s$ 得
\[
r_s=\frac{Y_{s,1}-Y_{s,2}-\sigma_T}{Y_{s,1}+Y_{s,2}+\sigma_T}.
\]
对 TM 波，用标准 $p$ 极化单位矢量定义入射和反射幅；反射时基矢的切向分量随传播方向约定改变，边界代数给
\[
r_p=\frac{Y_{p,2}-Y_{p,1}+\sigma_L}{Y_{p,1}+Y_{p,2}+\sigma_L}.
\]
令 $\sigma_{L,T}\to0$，两式退化为普通 Fresnel 系数；令分母为零，则恢复正文表面本征模条件。

平面平移对称使散射 Green 张量对面内坐标差作二维 Fourier 对角化。源到界面再反射回观察点的法向传播相位是
\[
\ee^{\ii k_{z1}z'}\ee^{\ii k_{z1}z}=\ee^{\ii k_{z1}(z+z')}.
\]
自由一维 Helmholtz 预解式给 $\ii/(2k_{z1})$，再按 TE/TM 完备分解便得到正文\eqnrefs{eq:2d_sheet_reflected_green_weyl}。

对匀速平行电子，
\[
\int_{-\infty}^{\infty}\dd t\,\ee^{\ii\omega t-\ii q_xvt}
=2\pi\delta(\omega-q_xv),
\]
所以 $q_x=\omega/v$。若 $k_{z1}=\ii\kappa_1$ 且 $z=z'=b$，则
\[
\left|\ee^{\ii k_{z1}(z+z')}\right|=\ee^{-2\kappa_1b},
\]
得到正文\eqnrefs{eq:2d_eels_evanscent_impact_parameter_factor}。
\end{derivation}

\begin{derivation}[极性晶格的纵横光学频率关系]
从正文\eqnrefs{eq:polar-optical-oscillator-eom-dp76} 的受驱极性光学模式开始：
\[
M_r(\ddot u+\Gamma_{\rm ph}\dot u+\omega_{\rm TO}^2u)=Z^*eE,
\qquad
P_{\rm ion}=n_{\rm ion}Z^*eu.
\]
采用全文的 $\ee^{-\ii\omega t}$ 约定，时间导数映射为
\[
\dot u\mapsto-\ii\omega u(\omega),
\qquad
\ddot u\mapsto-\omega^2u(\omega).
\]
因此运动方程变成代数方程
\[
M_r\left(\omega_{\rm TO}^2-\omega^2-\ii\Gamma_{\rm ph}\omega\right)u(\omega)
=Z^*eE(\omega),
\]
解得
\[
u(\omega)
=\frac{Z^*e/M_r}
{\omega_{\rm TO}^2-\omega^2-\ii\Gamma_{\rm ph}\omega}E(\omega).
\]
离子极化于是为
\begin{align*}
P_{\rm ion}(\omega)
&=n_{\rm ion}Z^*eu(\omega)\\
&=\frac{n_{\rm ion}Z^{*2}e^2/M_r}
{\omega_{\rm TO}^2-\omega^2-\ii\Gamma_{\rm ph}\omega}E(\omega).
\end{align*}

高频电子极化已经吸收到背景介电常数 $\varepsilon_\infty$ 中，所以位移场为
\[
D(\omega)=\varepsilon_0\varepsilon_\infty E(\omega)+P_{\rm ion}(\omega)
=\varepsilon_0\varepsilon(\omega)E(\omega).
\]
除以 $\varepsilon_0E$：
\[
\varepsilon(\omega)
=\varepsilon_\infty
+\frac{n_{\rm ion}Z^{*2}e^2/(\varepsilon_0M_r)}
{\omega_{\rm TO}^2-\omega^2-\ii\Gamma_{\rm ph}\omega}.
\]
通分后
\begin{align*}
\varepsilon(\omega)
=\varepsilon_\infty
\frac{\omega_{\rm TO}^2-\omega^2-\ii\Gamma_{\rm ph}\omega
+n_{\rm ion}Z^{*2}e^2/(\varepsilon_0\varepsilon_\infty M_r)}
{\omega_{\rm TO}^2-\omega^2-\ii\Gamma_{\rm ph}\omega}.
\end{align*}
定义长波纵向光学频率
\[
\omega_{\rm LO}^2
\equiv\omega_{\rm TO}^2
+\frac{n_{\rm ion}Z^{*2}e^2}
{\varepsilon_0\varepsilon_\infty M_r},
\]
分子立即化为
$\omega_{\rm LO}^2-\omega^2-\ii\Gamma_{\rm ph}\omega$，所以
\[
\varepsilon(\omega)
=\varepsilon_\infty
\frac{\omega_{\rm LO}^2-\omega^2-\ii\Gamma_{\rm ph}\omega}
{\omega_{\rm TO}^2-\omega^2-\ii\Gamma_{\rm ph}\omega},
\]
即正文\eqnrefs{eq:polar_optical_dielectric_response}。

无阻尼时，$\omega=\omega_{\rm TO}$ 使分母为零，对应横向光学共振；$\omega=\omega_{\rm LO}$ 使 $\varepsilon(\omega)=0$，对应无外加自由电荷时纵向条件
$D=\varepsilon_0\varepsilon E=0$ 的非零电场解。静态极限再取 $\omega=0$、$\Gamma_{\rm ph}=0$：
\[
\varepsilon_s\equiv\varepsilon(0)
=\varepsilon_\infty
\frac{\omega_{\rm LO}^2}{\omega_{\rm TO}^2}.
\]
整理得到
\[
\frac{\varepsilon_s}{\varepsilon_\infty}
=\frac{\omega_{\rm LO}^2}{\omega_{\rm TO}^2},
\]
即正文\eqnrefs{eq:lattice_lst_relation}。
\end{derivation}
